\documentclass{article}
\usepackage[left=2cm, right=2cm, top=2cm, bottom=2cm]{geometry}
\usepackage{graphicx}
\usepackage{amsmath}
\usepackage{amssymb}
\usepackage{amsthm}
\usepackage{fancyhdr}
\usepackage{verbatim}
\usepackage{listings}
\usepackage{xcolor}
\usepackage{pdfpages}
\usepackage{cancel}
\usepackage{physics}
\usepackage{esint}
\usepackage{braket}

\lstset{
    language=C++,
    basicstyle=\ttfamily\footnotesize,
    keywordstyle=\color{blue},
    commentstyle=\color{green},
    stringstyle=\color{red},
    numbers=left,
    numberstyle=\tiny\color{gray},
    frame=single,
    breaklines=true,
    captionpos=b,
}


\title{PHYS 427: Lecture \# 12}
\author{Cliff Sun}

\newtheorem{theorem}{Theorem}[section]
\newtheorem{lemma}[theorem]{Lemma}
\newtheorem{definition}[theorem]{Definition}
\newtheorem{conjecture}[theorem]{Conjecture}
\newtheorem{proposition}[theorem]{Proposition}
\newtheorem{corollary}[theorem]{Corollary}
\newtheorem{one minute paper}[theorem]{One Minute Paper}

\pagestyle{fancy}
\lhead{\textbf{\thepage}\ \ \nouppercase{\rightmark}}
\chead{PHYS 427: Lecture \# 12}
\rhead{Cliff Sun}

\begin{document}

\maketitle

\section*{Lecture Span}
\begin{enumerate}
    \item Grand Canonical Ensemble
\end{enumerate}

For midterm, $\Omega, S, T, Z, F$. Up to and including lecture 10. 

\section*{Grand Canonical Ensemble}

Reservoir and system are coupled. Note that $U_{\rm tot}$ and $N_{\rm tot}$ are fixed. Also, $S$ is in a single microstate $i$ with energy $\epsilon_i$, $N_i$ particles. 

The probability for $S$ to be in that state is 

\begin{equation}
    p(N_i, \epsilon_i) = \frac{\Omega_{R + S}(N_i, \epsilon_i)}{\sum \Omega_{R + S}}
\end{equation}

Moreover, 

\begin{equation}
    \Omega_{R + S} = \Omega_R \Omega_S
\end{equation}

Then, 

\begin{align}
    \frac{p(N_1, \epsilon_1)}{p(N_2, \epsilon_2)} &= \frac{\Omega_R(N_0 - N_1, U_0 - \epsilon_1)}{\Omega_R(N_0 - N_2, U_0 - \epsilon_2)} \\
    &= \frac{\exp(S_R(\dots)/k_B)}{\exp(S_R(\dots)/k_B)} \\
    &= \exp((S_R(N_1) - S_R(N_2))/k_B)
\end{align}

Then, expand about $U_0$ and $N_0$, then 

\begin{align}
    S_R &\approx S_R(U_0, N_0) - \frac{\partial S_R}{\partial N} \cdot N_1 - \frac{\partial S}{\partial U}\cdot \epsilon_1\\
    &= S_R + \frac{\mu}{T}N_1 - \frac{\epsilon_1}{T}
\end{align}

Then, 

\begin{equation}
    \frac{p_1}{p_2} = \frac{\exp\left( \frac{\mu N_1 - \epsilon_1}{k_B T} \right)}{\exp(\frac{\mu N_2 - \epsilon_2}{k_B T})}
\end{equation}

Moreover, we know that 

\begin{equation}
    \sum_i \frac{p_i}{p_2} = \frac{1}{p_2}
\end{equation}

Therefore, 

\begin{equation}
    p_2 = \frac{\exp(-(\epsilon_1 - \mu N_1)/k_B T)}{\mathcal{Z}}
\end{equation}

Here, $\mathcal{Z}$ is the Gibbs sum and defined as: 

\begin{equation}
    \mathcal{Z} = \sum_{N_i, \epsilon_i(N_i)}e^{-\left( \epsilon_i - \mu N_i \right)/k_B T}
\end{equation}

The top is called the Gibbs factor. Also, 

\begin{equation}
    \mathcal{Z} = \sum_N \exp(\mu N/k_B T) Z(N,T)
\end{equation}

Also,

\begin{equation}
    \langle N \rangle = \lambda \frac{\partial}{\partial \lambda} \ln (\mathcal{Z}(\lambda, T))
\end{equation}

Where 

\begin{equation}
    \lambda = \exp(\mu/k_B T)
\end{equation}

\begin{equation}
    \langle \epsilon \rangle = \frac{\mu}{\beta}\partial_{\mu} \ln \mathcal{Z} - \partial_{\beta}\ln \mathcal{Z}
\end{equation}

Next, we calculate 

\begin{enumerate}
    \item $\mathcal{Z}$ for myoglobin
    \item probability that a single $O_2$ bounds to it
\end{enumerate}

Note, the \# of $O_2$ binded to the myoglobin is either $N=0$ or $N=1$. Then, the gibbs sum is 

\begin{equation}
    \mathcal{Z} = 1 + \exp(-\beta(\epsilon_{O_2} - \mu))
\end{equation}

The prob. of 1 $O_2$ binding to the myoglobin is 

\begin{equation}
    f = \frac{1}{1 + \lambda^{-1}\exp(\epsilon_{O_2}/k_B T)}
\end{equation}

Then, assume that $O_2$ is in equilibrium with the Mb and $O_2$, then treat oxygen as an ideal gas and calculate 

\begin{equation}
    \mu(O_2) = k_B T \ln\left( n_{O_2}/n_Q \right)
\end{equation}

Then, 

\begin{equation}
    f = \frac{p}{p + n_Q k_B T \exp(\epsilon_{n_Q}/k_B T)}
\end{equation}

Here, $p$ is the partial pressure of Oxygen. 

\section*{Formalism}

Assume two chemical species $N^0_A$ and $N^0_B$, then 

\begin{equation}
    \mathcal{Z} = \sum_{N_A, N_B, \dots}\exp((\mu_A N_A + \dots)/k_B T)\sum_{\epsilon(N_A, N_B,\dots)}\exp(-\epsilon(N_A,N_B,\dots)/k_B T)
\end{equation}

\section*{Carbon monoxide poisoning}

Assume $O_2$ and $CO$ in the air. They are competing to attach to your myoglobin. But $CO_2$ has a lower binding energy. The probability of $CO_2$ binding to Mb is greater than that of $O_2$. 

\end{document}